\documentclass[a3paper,landscape,10pt]{article}
\usepackage[margin=10mm]{geometry}
\usepackage[french]{babel}
\usepackage{amsmath,amssymb}
\makeatletter
\def\input@path{{../../../Style/}}
\makeatother
\NeedsTeXFormat{LaTeX2e}
\ProvidesPackage{TLPosterCarteMentale}[2026/08/03 Posters et cartes mentales SI]

\RequirePackage{tikz}
\RequirePackage[most]{tcolorbox}
\RequirePackage{xcolor}
\RequirePackage{amsmath,amssymb}
\RequirePackage{graphicx}
\usetikzlibrary{positioning,arrows.meta,calc,shapes.geometric,mindmap,backgrounds,fit}

\definecolor{TLPosterBlue}{RGB}{24,86,141}
\definecolor{TLPosterLight}{RGB}{234,243,250}
\definecolor{TLPosterAccent}{RGB}{232,126,34}
\definecolor{TLPosterGreen}{RGB}{52,152,102}
\definecolor{TLPosterRed}{RGB}{192,57,43}
\definecolor{TLPosterGray}{RGB}{80,90,100}

\newcommand{\TLPosterTitre}[2]{%
  \begin{tcolorbox}[enhanced,boxrule=0pt,arc=0pt,colback=TLPosterBlue,left=6mm,right=6mm,top=4mm,bottom=4mm]
    {\color{white}\bfseries\LARGE #1}\hfill{\color{white}\large #2}
  \end{tcolorbox}%
}

\newtcolorbox{TLPosterBloc}[2][]{
  enhanced,breakable,
  colback=TLPosterLight,
  colframe=TLPosterBlue,
  colbacktitle=TLPosterBlue,
  coltitle=white,
  fonttitle=\bfseries,
  title={#2},
  boxrule=0.8pt,
  arc=2mm,
  left=3mm,right=3mm,top=2mm,bottom=2mm,
  #1
}

\newtcolorbox{TLPosterFormule}[1][]{
  enhanced,
  colback=white,
  colframe=TLPosterAccent,
  boxrule=1pt,
  arc=2mm,
  fontupper=\bfseries,
  halign=center,
  #1
}

\newcommand{\TLPosterMethode}[1]{%
  \begin{tcolorbox}[enhanced,colback=TLPosterGreen!8,colframe=TLPosterGreen,title=Méthode,fonttitle=\bfseries]
  #1
  \end{tcolorbox}%
}

\newcommand{\TLPosterAttention}[1]{%
  \begin{tcolorbox}[enhanced,colback=TLPosterRed!6,colframe=TLPosterRed,title=Point de vigilance,fonttitle=\bfseries]
  #1
  \end{tcolorbox}%
}

\tikzset{
  TLmindroot/.style={concept,concept color=TLPosterBlue,text=white,font=\bfseries\large,minimum size=34mm},
  TLmindlevel1/.style={level 1 concept/.append style={font=\bfseries,minimum size=23mm,sibling angle=45}},
  TLmindlevel2/.style={level 2 concept/.append style={font=\small,minimum size=17mm}},
  TLarrow/.style={-{Stealth[length=3mm]},thick,draw=TLPosterBlue},
  TLbox/.style={rounded corners=2mm,draw=TLPosterBlue,fill=TLPosterLight,align=center,inner sep=3mm}
}

\newenvironment{TLCarteMentale}[1]{%
  \begin{tikzpicture}[mindmap,grow cyclic,concept color=TLPosterBlue,every node/.style={concept},TLmindlevel1,TLmindlevel2]
  \node[TLmindroot]{#1}
}{;
  \end{tikzpicture}
}

\newcommand{\TLBranche}[3][]{child[concept color=#2]{node[#1]{#3}}}

\endinput

\IfFileExists{TLLogiqueSED.sty}{\NeedsTeXFormat{LaTeX2e}
\ProvidesPackage{TLLogiqueSED}[2026/08/03 Logique combinatoire et SED]
\RequirePackage{tikz}
\usetikzlibrary{arrows.meta,positioning,calc,shapes.geometric,fit}

\newcommand{\TLPortesLogiques}{%
\begin{tikzpicture}[>=Latex,node distance=13mm]
\node[draw,rounded corners,minimum width=18mm,minimum height=9mm] (and) {ET};
\node[draw,rounded corners,right=of and,minimum width=18mm,minimum height=9mm] (or) {OU};
\node[draw,rounded corners,right=of or,minimum width=18mm,minimum height=9mm] (not) {NON};
\node[left=8mm of and] (a) {$a,b$}; \draw[->] (a)--(and);
\node[below=4mm of and] {$s=a\cdot b$};
\node[below=4mm of or] {$s=a+b$};
\node[below=4mm of not] {$s=\overline a$};
\end{tikzpicture}}

\newcommand{\TLDemarcheCombinatoire}{%
\begin{tikzpicture}[>=Latex,node distance=8mm]
\foreach \n/\t in {1/Cahier des charges,2/Table de vérité,3/Équation logique,4/Simplification,5/Implantation}{
\node[draw,rounded corners,minimum width=32mm,minimum height=8mm] (n\n) at (0,-1.15*\n) {\t};}
\foreach \a/\b in {1/2,2/3,3/4,4/5}{\draw[->] (n\a)--(n\b);}
\end{tikzpicture}}

\newcommand{\TLGrapheEtatsSimple}{%
\begin{tikzpicture}[>=Latex,node distance=30mm]
\node[draw,circle,minimum size=16mm] (e0) {Attente};
\node[draw,circle,right=of e0,minimum size=16mm] (e1) {Marche};
\draw[->] ([xshift=-13mm]e0.west)--(e0.west);
\draw[->,bend left=18] (e0) to node[above]{départ} (e1);
\draw[->,bend left=18] (e1) to node[below]{arrêt} (e0);
\end{tikzpicture}}

\newcommand{\TLDiagrammeSequenceSimple}{%
\begin{tikzpicture}[>=Latex]
\foreach \x/\t in {0/Utilisateur,3/Commande,6/Actionneur}{
\node[draw,minimum width=22mm] (h\x) at (\x,0) {\t};
\draw[dashed] (\x,-.35)--(\x,-4);}
\draw[->] (0,-1)--node[above]{ordre}(3,-1);
\draw[->] (3,-2)--node[above]{commander}(6,-2);
\draw[->,dashed] (6,-3)--node[above]{retour}(3,-3);
\end{tikzpicture}}

\newcommand{\TLCodeurIncremental}{%
\begin{tikzpicture}[>=Latex]
\draw[thick] (0,0) circle (1.25);
\foreach \a in {0,30,...,330}{\draw[thick] (0,0)--(\a:1.25);}
\draw[fill=white] (0,0) circle (.45);
\node[right] at (1.6,.35) {voie A};
\node[right] at (1.6,-.35) {voie B};
\draw[->] (1.45,.35)--(1,.35);
\draw[->] (1.45,-.35)--(1,-.35);
\end{tikzpicture}}

\newcommand{\TLModeleOSI}{%
\begin{tikzpicture}
\foreach \i/\t in {1/Physique,2/Liaison,3/Réseau,4/Transport,5/Session,6/Présentation,7/Application}{
\node[draw,minimum width=36mm,minimum height=6.5mm] at (0,.7*\i) {\i\ -- \t};}
\end{tikzpicture}}
}{}
\pagestyle{empty}
\begin{document}
\TLPosterTitre{Logique et systèmes séquentiels}{Fiche résumé éditable}
\begin{multicols}{2}

\begin{TLPosterBloc}{Variables logiques}
Une variable binaire prend les valeurs $0$ ou $1$. Une fonction logique associe une sortie binaire à une combinaison des entrées.
\end{TLPosterBloc}

\begin{TLPosterBloc}{Opérateurs fondamentaux}
NON : $\overline{a}$ ; ET : $a\cdot b$ ; OU : $a+b$ ; OU exclusif : $a\oplus b$. Les portes NAND et NOR sont universelles.
\end{TLPosterBloc}

\begin{TLPosterBloc}{Algèbre de Boole}
Identités utiles : $a+a=a$, $a\cdot a=a$, $a+\overline a=1$, $a\cdot\overline a=0$ et absorption $a+a b=a$.
\end{TLPosterBloc}

\begin{TLPosterBloc}{Lois de De Morgan}
\[
\overline{a+b}=\overline a\,\overline b,
\qquad
\overline{a b}=\overline a+\overline b.
\]
Elles permettent de transformer une réalisation logique et de basculer entre portes ET, OU, NAND et NOR.
\end{TLPosterBloc}

\begin{TLPosterBloc}{Table de vérité}
Pour $n$ entrées, la table comporte $2^n$ combinaisons. Elle permet de passer du cahier des charges à l'équation logique puis au schéma de portes.
\end{TLPosterBloc}

\begin{TLPosterBloc}{Formes canoniques}
La somme de produits est construite à partir des mintermes associés aux lignes où la sortie vaut $1$. Le produit de sommes utilise les maxtermes.
\end{TLPosterBloc}

\begin{TLPosterBloc}{Simplification}
Utiliser les identités de Boole ou un tableau de Karnaugh. Regrouper les cases adjacentes par paquets de $1$, $2$, $4$, $8$, etc., aussi grands que possible.
\end{TLPosterBloc}

\begin{TLPosterBloc}{Logique combinatoire}
La sortie dépend uniquement des entrées présentes. Exemples : multiplexeur, décodeur, additionneur, comparateur et codeur.
\end{TLPosterBloc}

\begin{TLPosterBloc}{Mémoire et logique séquentielle}
La sortie dépend des entrées et de l'état précédent. Les bascules RS, D, JK ou T mémorisent un bit et constituent registres, compteurs et mémoires.
\end{TLPosterBloc}

\begin{TLPosterBloc}{Horloge et chronogramme}
Les changements d'état peuvent être synchronisés sur un front montant ou descendant. Respecter les temps de propagation, de préparation et de maintien.
\end{TLPosterBloc}

\begin{TLPosterBloc}{Automate fini}
Un automate est défini par ses états, son état initial, ses transitions et ses sorties. Une transition peut porter la forme : événement $[\text{garde}]$ / effet.
\end{TLPosterBloc}

\begin{TLPosterBloc}{Moore et Mealy}
Dans une machine de Moore, les sorties dépendent uniquement de l'état. Dans une machine de Mealy, elles dépendent de l'état et des entrées.
\end{TLPosterBloc}

\begin{TLPosterBloc}{Diagramme d'états SysML}
Identifier les états stables, les événements déclencheurs, les conditions de garde, les actions d'entrée, de sortie et les effets des transitions.
\end{TLPosterBloc}

\begin{TLPosterBloc}{Diagramme de séquence}
Il représente les échanges de messages entre acteurs et blocs au cours du temps : lignes de vie, messages synchrones ou asynchrones, fragments alternatifs et boucles.
\end{TLPosterBloc}

\begin{TLPosterBloc}{Communication numérique}
Distinguer information, codage, support et protocole. Les transmissions peuvent être série ou parallèle, synchrones ou asynchrones, simplex, semi-duplex ou duplex.
\end{TLPosterBloc}

\begin{TLPosterBloc}{Démarche de résolution}
Lister entrées et sorties, établir la table de vérité ou les états, simplifier, réaliser le schéma ou l'automate, tracer les chronogrammes puis vérifier tous les cas limites.
\end{TLPosterBloc}

\end{multicols}
\end{document}
